\documentclass[11pt]{article}
\usepackage[margin=1in]{geometry}
\usepackage{amsmath,amssymb}
\usepackage{booktabs}
\usepackage{graphicx}
\usepackage{enumitem}
\usepackage{setspace}
\usepackage[hidelinks]{hyperref}
\usepackage{url}
\usepackage{xurl}
\usepackage{caption}
\usepackage{subcaption}
\usepackage{threeparttable}
\usepackage{tabularx}
\usepackage{float}
\graphicspath{{../input/analysis_figures/}{../input/static_figures/}{../input/tables/}{../input/}}

\setlength{\parskip}{6pt}
\setlength{\parindent}{0pt}
\onehalfspacing

% Convenience macros
\newcommand{\E}{\mathbb{E}}
\newcommand{\R}{\mathbb{R}}

\begin{document}

\begin{center}
{\LARGE\bfseries Markups and Public Procurement:\\[4pt]
Evidence from Czech Construction Tenders}\\[16pt]
{\large Marek Chadim\footnote{Pre-Doctoral Fellow, Tobin Center for Economic Policy, Yale University. Email: \texttt{marek.chadim@yale.edu}. I thank Chad Syverson for invaluable guidance, Steven Berry and Jesse Shapiro for comments, and participants in ECON 601 (Graduate IO II) at Yale. Replication package: \url{https://github.com/marek-chadim/Markups-and-Public-Procurement}.}}\\[8pt]
{\normalsize Yale University}\\[12pt]
{\normalsize March 2026}\\[8pt]
{\normalsize\textsc{Preliminary and Incomplete}}
\end{center}

\bigskip

\begin{abstract}
\noindent
This paper estimates firm-level markups in the Czech construction sector and measures the causal effect of public procurement participation on markups. Combining firm-level financial statements with the universe of Czech public procurement contracts, I estimate markups via the production function approach of De Loecker and Warzynski (2012) under the ACF (2015) control function framework. The primary specification---firm fixed effects on a panel of 1,267 firms with 6,849 firm-year observations---yields a procurement markup premium of approximately 14\%, with Oster (2019) bounds confirming robustness to omitted variable bias. Following De Loecker (2013) and De Loecker and Syverson (2021, \emph{Handbook of Industrial Organization}), I include the procurement indicator in the Markov productivity transition---the specification required for internal consistency whenever the treatment may affect productivity dynamics through learning-by-doing, scale, or technology transfer. Under this specification, the regression-adjusted premium is 14\%. A specification that omits procurement from the law of motion imposes the strong assumption that procurement has no effect on productivity, and is internally inconsistent with any hypothesis of learning-by-procuring. The premium declined from 30\% in 2006 to 10\% in 2021, with DiD estimates of $-2.6$ and $-2.2$ log points associated with Act 55/2012 (single-bid ban) and Act 134/2016 (MEAT criteria) respectively. I complement the firm-level analysis with aggregate sales-weighted markup trends constructed following De Loecker, Eeckhout, and Unger (2020), decomposing the evolution into within-firm, between-firm reallocation, and net-entry components. A systematic specification sensitivity analysis reveals that the variable input definition and the Markov process specification dominate all other modeling choices.

\medskip
\noindent \textbf{JEL Codes:} D24, H57, J42, L13, L74

\noindent \textbf{Keywords:} markups, production function estimation, public procurement, Czech Republic, construction
\end{abstract}

\newpage
\tableofcontents
\newpage

%% ============================================================
%% 1. INTRODUCTION
%% ============================================================
\section{Introduction}\label{sec:intro}

Do rising markups reflect growing market power? De Loecker, Eeckhout, and Unger \cite{DeLoekerEeckhout2020} document that US markups rose from 21\% above marginal cost in 1980 to 61\% in 2016, concentrated among high-markup firms. Benkard, Miller, and Yurukoglu \cite{BMY2026} challenge this finding: the estimated trend is sensitive to production function specification choices---particularly the variable input definition and overhead cost treatment. In their reply, De Loecker, Eeckhout, and Unger \cite{DLEU2025reply} trace the divergence to outliers in a single four-digit industry and show that output elasticities are stable over time, with the trend driven entirely by expenditure shares. The exchange reveals a deeper question: if markup \emph{levels} are specification-sensitive, can markup-based \emph{treatment effects} be trusted?

This paper brings the production approach to a setting where that question can be answered directly. Czech construction procurement provides a binary treatment---procurement participation---against which to measure whether specification sensitivity in markup estimation propagates to causal estimates. Three findings emerge. First, firms participating in public procurement charge markups approximately 14\% higher than non-participants, with an event study showing flat pre-trends and a sharp jump at entry. Second, this premium is robust across nine distinct estimation methods: while markup \emph{levels} range from 1.02 to 2.17, the treatment effect is stable at 12.5--14.3\% across all methods except OLS. Third, the premium has declined from 30\% in 2006 to 10\% in 2021, tracking a sequence of Czech transparency reforms.

Czech construction is a clean testbed for the production approach. Cost of goods sold---dominated by cement, steel, concrete, and fuel---constitutes approximately 50\% of gross output, providing a well-measured variable input with stable cost shares. The industry has been studied extensively in the markup literature: Collard-Wexler and De Loecker \cite{CWDLKvsOmega} analyze US steel, while Kroft, Luo, Mogstad, and Setzler \cite{KLMS2025} estimate a double market power model for US construction procurement. Between 2006 and 2021, Czech policymakers enacted a series of transparency reforms---the 2012 single-bidding ban, the 2016 shift to MEAT criteria, and the 2017 Register of Contracts---providing time-series variation in the scope for discretion and favoritism.

The estimation follows De Loecker and Warzynski \cite{DLW2012}: the markup equals the ratio of the output elasticity of a variable input to its expenditure share. Output elasticities are recovered from a translog production function estimated via Ackerberg, Caves, and Frazer \cite{ACF2015}, with the productivity law of motion conditioned on lagged procurement status following De Loecker \cite{DeLoecker2013}. The specification sensitivity analysis---varying the variable input, Markov process, functional form, and estimator---reveals that the premium is identified from cross-firm variation in cost shares ($\alpha^V_{it}$, observed data), not from the estimated output elasticity ($\theta^V$). This is consistent with the Gandhi, Navarro, and Rivers \cite{GNR2020} non-identification result and explains why the treatment effect survives even when markup levels are specification-sensitive.

This paper contributes to the literature on markup measurement, treatment evaluation with production function methods, and Czech procurement policy. On measurement, the nine-method comparison extends the Benkard, Miller, and Yurukoglu \cite{BMY2026} specification sensitivity analysis from aggregate trends to treatment effects, showing that the latter are more robust than the former. On treatment evaluation, the DLW-style entry/exit decomposition \cite{DLW2012} and productivity accounting extend their export analysis to procurement, documenting that productivity explains 76\% of the markup association (closely matching their 70\% for exports). Conlon, Miller, Otgon, and Yao \cite{CMOY2023} show that production-approach markup trends are uncorrelated with price changes ($R^2 < 0.001$), raising questions about what markups measure; the treatment evaluation design sidesteps this concern by focusing on within-firm changes around a discrete event. On policy, the declining premium provides evidence that transparency reforms reduce the scope for rent extraction in procurement, complementing Baran\v{e}k and Titl's \cite{BaranekTitl2024} finding that political connections raise Czech procurement prices by 6\%. Unlike De Loecker and Scott \cite{DeLoeckerScott2025}, who combine production and demand data for the US brewing industry, the bid-level data here include only winning prices, limiting the analysis to the production approach---making specification sensitivity a first-order concern.


%% ============================================================
%% 2. SETTING AND DATA
%% ============================================================
\section{Setting and Data}\label{sec:data}

\subsection{Czech Construction}\label{sec:setting}

The Czech construction sector accounts for 49.6\% of all public procurement contracts, approximately 7\% of GDP, and employs over 7\% of the labor force. It is also among the least competitive procurement markets in the European Union: single-bidder tenders exceed 48\% of awarded contracts, and each additional bidder reduces prices by roughly 3.3\% \cite{PalgutaPertold2017}. The sector comprises three CZ-NACE sub-industries: NACE 41 (construction of buildings), NACE 42 (civil engineering), and NACE 43 (specialized construction activities).

Between 2006 and 2021, Czech policymakers enacted a sequence of transparency reforms aimed at increasing competition and reducing the scope for favoritism. Table \ref{tab:reforms} summarizes the key legislative changes.

\begin{table}[htbp]\centering
\caption{Czech Procurement Transparency Reforms, 2006--2021}\label{tab:reforms}
\begin{threeparttable}
\begin{tabular}{llp{8cm}}
\toprule
Date & Reform & Key Change \\
\midrule
1 Jul 2006 & Act 137/2006 & Full EU procurement framework; TED publication requirement \\
\textbf{1 Apr 2012} & \textbf{Act 55/2012} & \textbf{Mandatory cancellation of single-bid tenders}; 6--10\% price reduction (Titl 2023) \\
1 Jan 2014 & Act 134/2013 & Bearer share abolition; anonymous ownership eliminated \\
1 Oct 2016 & Act 134/2016 & Shift from lowest-price toward most economically advantageous tender (MEAT) criteria \\
\textbf{1 Jul 2017} & \textbf{Register of Contracts} & \textbf{Contract validity contingent on online publication} \\
1 Jan 2018 & Beneficial Ownership Register & Public register of ultimate beneficial owners \\
\bottomrule
\end{tabular}
\begin{tablenotes}\footnotesize
\item \emph{Notes:} Bold entries denote reforms most relevant to the declining premium. The 2012 single-bidding ban was the most impactful: Titl \cite{Titl2023} estimates a 6--10\% reduction in procurement prices. The 2017 Register of Contracts made contract publication a validity condition.
\end{tablenotes}
\end{threeparttable}
\end{table}

These reforms provide exogenous time-series variation in competition intensity. Palguta and Pertold-Gebicka \cite{PalgutaPertold2017} document that below-threshold manipulation allowed politically connected firms to avoid competitive tendering before the reforms, while Baran\v{e}k and Titl \cite{BaranekTitl2024} show that political connections raise procurement prices by 6\%, with EU co-funding oversight fully mitigating this overpricing.

\subsection{Data Sources}\label{sec:datasources}

Three Czech data sources are combined, covering 2005--2023.

\textbf{Firm financials.} The MagnusWeb database provides annual financial statements for Czech firms classified in CZ-NACE division F (construction). Key variables include sales (GO), cost of goods sold (COGS), number of employees (L), and fixed assets (K). The raw dataset also contains wages (W) and intermediate inputs (II) separately at full coverage, enabling decomposition of GO $\approx$ W + II. All nominal variables are deflated using OECD industry deflators.

\textbf{Public procurement.} The Datlab register contains the universe of Czech public procurement contracts, compiled from the official ISVZ system. For each tender, the data record the contract value, signature date, procurement procedure, awarding entity, and the identities of all bidders. Procurement records are merged to firm financials using company identifiers. The primary treatment variable is $pp_{it}$, a binary indicator equal to one if firm $i$ has an active procurement contract in year $t$.

\textbf{Sample construction.} The estimation sample is rebuilt from raw CSVs using a transparent Python pipeline (\texttt{rebuild\_data.py}). Key steps include: (i) 1:1 merges with deduplication, (ii) tender year assigned from signature date, (iii) winsorization at the 2nd and 98th percentiles by NACE 2-digit on log variables following De Ridder, Grassi, and Morzenti \cite{DGM2026}, and (iv) trimming observations with overhead-to-gross-output ratios exceeding 0.90. The resulting sample contains 1,521 firms and 9,164 firm-year observations.

\textbf{Orbis financial statements.} To overcome MagnusWeb's limitation of categorical employment brackets, I supplement with Bureau van Dijk's Orbis database accessed through WRDS. Orbis provides continuous employee headcounts, cost of employees, operating revenue, cost of goods sold, and fixed assets for Czech firms across all industries. Merging on company identifiers (I\v{C}O), 84\% of MagnusWeb construction firms appear in Orbis (1,260 of 1,498). Capital correlates at 0.983 and wages at 0.951 between the two sources.

Table \ref{tab:summary} reports summary statistics for the estimation sample. Panel~B compares treated and control firms: procurement participants are larger but operate at similar cost structures, consistent with selection on scale rather than efficiency.

\input{../input/tables/summary_stats.tex}

\subsection{Data Limitations}\label{sec:limitations}

Four limitations bear mention. First, the Datlab tender data include only winning bid prices, not losing bids, precluding estimation of a full demand system or structural auction model and restricting the analysis to the production approach of De Loecker and Warzynski \cite{DLW2012}. Second, MagnusWeb employment is reported in categorical brackets rather than continuous counts, limiting labor supply elasticity estimation. Third, the tender data do not flag EU co-funding, precluding a direct test of the intergovernmental oversight channel identified by Baran\'{e}k and Titl \cite{BaranekTitl2024}. Fourth, output is measured as revenue, not physical quantity. Following De Ridder, Grassi, and Morzenti \cite{DGM2026}, the revenue-to-quantity bias is a sector-level constant that shifts markup levels but preserves cross-firm variation---the procurement premium as a ratio is therefore robust to this bias.


%% ============================================================
%% 3. PRODUCTION APPROACH TO MARKUP ESTIMATION
%% ============================================================
\section{Production Approach to Markup Estimation}\label{sec:methodology}

\subsection{Markup Formula}\label{sec:markup_formula}

Under cost minimization, a firm equates the output elasticity of each freely adjustable input to the product of its markup and the input's expenditure share. This yields the markup formula of De Loecker and Warzynski \cite{DLW2012}, building on Hall \cite{Hall1988}:
\begin{equation}\label{eq:markup}
    \mu_{it} = \frac{\theta^V_{it}}{\alpha^V_{it}},
\end{equation}
where $\theta^V_{it} = \partial \ln Y_{it} / \partial \ln V_{it}$ is the output elasticity of the variable input and $\alpha^V_{it} = P^V_{it} V_{it} / (P_{it} Y_{it})$ is the expenditure share of the variable input in total revenue. Under perfect competition $\mu_{it} = 1$ and $\theta^V_{it} = \alpha^V_{it}$; deviations from unity indicate market power. Appendix~\ref{app:markup_derivation} provides the full derivation from profit maximization.

The formula rests on three assumptions: (i) firms minimize variable cost conditional on output and fixed inputs; (ii) productivity is Hicks-neutral; and (iii) productivity evolves as a first-order Markov process conditional on lagged treatment status. The third assumption, established by De Loecker \cite{DeLoecker2013} for learning-by-exporting and extended by De Loecker and Syverson \cite{DLS2021} to general treatment evaluation, ensures internal consistency whenever the treatment may affect productivity dynamics.

\subsection{Estimation}\label{sec:estimation}

Estimating the output elasticity $\theta^V_{it}$ from production data confronts a fundamental simultaneity problem: firms observe their productivity $\omega_{it}$ and adjust variable inputs accordingly, biasing OLS upward for the variable input coefficient. Fixed effects do not solve this---they require the implausible assumption that productivity is time-invariant ($\omega_{it} = \omega_i$), and first-differencing exacerbates measurement error in capital, compressing $\hat{\beta}_k$ toward zero. Olley and Pakes \cite{OlleyPakes1996} propose a structural solution: invert a monotone policy function (investment) to recover $\omega_{it} = f_t^{-1}(k_{it}, i_{it})$, then use the Markov property of productivity to form moment conditions that identify the capital coefficient. Ackerberg, Caves, and Frazer \cite{ACF2015} refine this approach, replacing investment with the variable input demand function as the proxy and resolving a collinearity problem in the first stage. The Cobb-Douglas specification is:
\begin{equation}\label{eq:cd}
    y_{it} = \beta_0 + \beta_k k_{it} + \beta_{\text{cogs}} \text{cogs}_{it} + \omega_{it} + \varepsilon_{it},
\end{equation}
and the translog specification:
\begin{equation}\label{eq:tl}
    y_{it} = \beta_0 + \beta_k k_{it} + \beta_v v_{it} + \beta_{kk} k_{it}^2 + \beta_{vv} v_{it}^2 + \beta_{kv} k_{it} v_{it} + \omega_{it} + \varepsilon_{it},
\end{equation}
where $y_{it}$ is log deflated sales, $k_{it}$ is log capital, and $v_{it}$ is log variable input.

Estimation proceeds in two stages. The first stage recovers $\hat{\phi}_{it}$ from a flexible polynomial in inputs and state variables. The second stage forms moment conditions $\E[\xi_{it}(\beta) \cdot \mathbf{Z}_{it}] = 0$, where $\xi_{it} = \omega_{it} - g(\omega_{it-1}, pp_{it-1})$ is the innovation to the productivity Markov process and $\mathbf{Z}_{it} = (1, k_{it}, \text{cogs}_{it-1})$. The Markov transition includes the lagged procurement indicator:
\begin{equation}
    \hat{\omega}_{it} = \rho \hat{\omega}_{it-1} + \gamma pp_{it-1} + \delta \hat{p}_{it-1} + \xi_{it},
\end{equation}
where $\hat{p}_{it-1}$ is the predicted survival probability from a probit model, addressing the Olley-Pakes selection problem: surviving firms have systematically higher productivity innovations ($E[\xi_{it} | k_{it}, \text{active at } t] \neq 0$), biasing the Markov regression unless one conditions on the exit threshold \cite{OlleyPakes1996}. Including $pp_{it-1}$ in the law of motion is not a robustness check---it is the only specification consistent with the model under any hypothesis where procurement can affect productivity through learning-by-doing, scale economies, or technology transfer \cite{DeLoecker2013,DLS2021}. The control function inversion requires a scalar unobservable assumption: $\omega_{it}$ must be the only serially correlated component that the firm observes when choosing inputs. This rules out unobserved cross-firm input price variation and measurement error in the proxy variable---assumptions that are more plausible for materials (COGS) than for labor \cite{DGM2026}. Analytical standard errors follow the GMM sandwich formula of Ackerberg, Chen, and Hahn \cite{ACH2012}, with firm-clustered meat matrix and delta method for markup SEs.

\subsection{Variable Input Choice}\label{sec:variable_input}

The choice of variable input is the most consequential specification decision. Cost of goods sold (COGS) constitutes approximately 50\% of gross output in Czech construction, dominated by cement, steel, concrete, and fuel---materials plausibly under short-run adjustment. Intermediate inputs (II), at approximately 83\% of gross output, compress the expenditure share toward the output elasticity, mechanically pushing markups toward unity.

Raval \cite{Raval2023} formalizes the concern: under cost minimization with Hicks-neutral productivity, \emph{any} flexible input should recover the same markup. His finding that labor and materials markups are negatively correlated across seven countries rejects this overidentifying restriction. In the Czech construction panel, labor-based and COGS-based markups are positively correlated ($\hat{\beta} = 0.28$--$0.82$), preserving firm rankings even though $\beta = 1$ is rejected. The positive Czech correlation means the treatment effect is preserved even when overidentification fails.

De Ridder, Grassi, and Morzenti \cite{DGM2026} provide a further justification: if the variable input faces a monopsony wedge, the production approach measures the product of the markup and the input wedge. Labor markets may feature such wedges \cite{KLMS2025}, making COGS the cleaner measure. The sensitivity analysis in Section~\ref{sec:nine_methods} confirms: using II instead of COGS flips the premium sign, while all other specification choices leave it intact.

\subsection{Specification Sensitivity}\label{sec:sensitivity}

The central methodological finding is that markup \emph{levels} are specification-sensitive but the procurement \emph{premium} is not. The sensitivity has a clean explanation. Under the Gandhi, Navarro, and Rivers \cite{GNR2020} non-identification result, $\theta^V$ is pinned down by the expenditure share $\alpha^V_{it}$, which is observed data. The premium is therefore identified from cross-firm variation in cost shares, not from the estimated production function. This identification from $\alpha$ is robust to the estimator (ACF, OLS, Blundell-Bond), the functional form (Cobb-Douglas, translog), and the imperfect competition correction \cite{ADL2024}. It fails only when an alternative variable input yields different cost shares---confirming that the variable input definition, not the estimator, is the binding specification decision.

De Loecker, Eeckhout, and Unger \cite{DLEU2025reply} make the same point for aggregate trends: output elasticities are stable, and all variation flows through the expenditure share. The Conlon, Miller, Otgon, and Yao \cite{CMOY2023} finding that markup growth is uncorrelated with PPI changes ($R^2 < 0.001$) underscores that cross-sectional cost-share variation, rather than production function estimation, drives markup measurement.

The Markov specification is the second key choice. Including the procurement dummy triples the regression-adjusted premium (from 0.003 to 0.138), because the plain specification attributes all procurement-induced productivity gains to the residual, biasing the markup premium toward zero. De Loecker and Syverson \cite{DLS2021} make this point with particular force: ``if the changes in the environment are part of the information set of a producer, they should be included in the model for forecasting productivity'' (p.~196).


%% ============================================================
%% 4. MARKUPS IN CZECH CONSTRUCTION
%% ============================================================
\section{Markups in Czech Construction}\label{sec:results}

\subsection{Production Function Estimates}\label{sec:pf_coefs}

Table \ref{tab:pf_estimates} reports the estimated production function coefficients. The baseline COGS elasticity ($\hat{\beta}_{\text{cogs}}$) is 0.957 for NACE 41, 1.100 for NACE 42, and 0.927 for NACE 43. Capital elasticities are small (0.044, $-$0.024, and 0.065 respectively). Returns to scale are close to unity (0.995--1.076), consistent with constant returns. First-stage $R^2$ values range from 0.89 to 0.95.

The OLS COGS coefficient is virtually identical to the ACF estimate: 0.973 versus 0.957 for NACE 41 ($\Delta = 1.7\%$), 1.108 versus 1.100 for NACE 42 ($\Delta = 0.7\%$). This near-equivalence is predicted by the Gandhi, Navarro, and Rivers \cite{GNR2020} non-identification theorem: proxy variable methods generically fail to identify the flexible input elasticity for gross output production functions. The translog specification yields similar mean elasticities but with substantially larger standard errors, reflecting weak identification of higher-order terms in the just-identified system. Following De Ridder, Grassi, and Morzenti \cite{DGM2026}, the Cobb-Douglas serves as the preferred specification.

\begin{table}[htbp]\centering
\caption{Production Function Estimates by Industry}\label{tab:pf_estimates}
\begin{threeparttable}
\begin{tabular}{l*{3}{c}}
\toprule
 & NACE 41 & NACE 42 & NACE 43 \\
 & (Buildings) & (Civil Eng.) & (Specialized) \\
\midrule
\multicolumn{4}{l}{\textit{Panel A: Baseline (survival + pp in Markov)}} \\
$\hat{\beta}_k$ & 0.044 & $-$0.024 & 0.065 \\
 & (0.003) & (0.014) & (0.005) \\
$\hat{\beta}_{\text{cogs}}$ & 0.957 & 1.100 & 0.927 \\
 & (0.007) & (0.017) & (0.015) \\
RTS & 1.001 & 1.076 & 0.992 \\
$N$ & 4,131 & 711 & 2,824 \\
\midrule
\multicolumn{4}{l}{\textit{Panel B: Robustness --- $\hat{\beta}_{\text{cogs}}$ (SE)}} \\
No survival & 0.969 (0.007) & 1.097 (0.017) & 0.926 (0.012) \\
No pp in Markov & 0.948 (0.003) & 1.104 (0.013) & 0.929 (0.014) \\
Plain ACF & 0.957 (0.003) & 1.102 (0.013) & 0.925 (0.012) \\
OLS & 0.973 & 1.108 & 0.899 \\
Translog & 1.459 (0.192) & 0.988 (0.212) & 1.181 (0.665) \\
\bottomrule
\end{tabular}
\begin{tablenotes}\footnotesize
\item \emph{Notes:} Two-step ACF estimator (Ackerberg et al.\ 2015). Baseline Markov: $\omega_t = \rho\omega_{t-1} + \gamma pp_{t-1} + \delta\hat{p}_{t-1} + \xi_t$. Instruments: $(1, k_t, \text{cogs}_{t-1})$. Analytical SEs (ACH 2012), clustered by firm. Translog reports mean output elasticity at sample midpoint.
\end{tablenotes}
\end{threeparttable}
\end{table}

\subsection{Markup Distributions}\label{sec:markups}

Average markups differ substantially across sub-industries: NACE 41 (buildings) averages 2.14, NACE 42 (civil engineering) 2.35, and NACE 43 (specialized construction) 1.71. These levels---prices 70--110\% above variable cost---may appear large for a competitive industry with 3--8\% profit margins. However, markups measure price relative to \emph{variable} cost only; fixed costs (land, equipment, overhead) drive the wedge. As Hall~\cite{Hall2018} emphasizes, an industry with 5\% profit margins can sustain markups of 2.0 if fixed costs account for 47\% of revenue.

The unconditional log markup difference between procurement and non-procurement firms is 0.186 (SE $= 0.003$). Controlling for year$\times$NACE fixed effects and input levels, the cross-sectional premium is $\hat{\beta}_1 = 0.140$ (SE $= 0.005$), implying a 14.9\% procurement markup premium. Czech construction markups fall within the range documented in the broader European literature \cite{DiezLeighTambunlertchai2018}.

\iffalse % BMY figure disabled --- depends on disabled pipeline script
\begin{figure}[htbp]\centering
\includegraphics[width=0.85\textwidth]{bmy_percentiles.pdf}
\caption{Evolution of Markups in Czech Construction, 2006--2021}\label{fig:markup_evolution}
\begin{minipage}{0.85\textwidth}\footnotesize
\emph{Notes:} Mean markup by year, computed from the baseline ACF Cobb-Douglas specification with survival correction and procurement dummy in the Markov transition. Shaded area: interquartile range.
\end{minipage}
\end{figure}
\fi

\subsection{Nine-Method Comparison}\label{sec:nine_methods}

Table \ref{tab:dls} confirms the robustness of the procurement premium across nine distinct markup estimation methods following the taxonomy of De Loecker and Syverson \cite{DLS2021}. The premium ranges from 12.5\% to 14.3\% across eight methods, with only OLS (0.3\%) as an outlier---consistent with the GNR non-identification result destroying cross-sectional variation under simultaneity bias. While markup \emph{levels} vary from 1.02 (cost-share) to 2.17 (Blundell-Bond), the treatment effect is stable. This validates the De Ridder, Grassi, and Morzenti \cite{DGM2026} prediction: revenue-based markup \emph{variation} is well-measured even when levels are biased by a sector-level constant.

\input{../input/tables/dls_comparison.tex}

\subsection{Aggregate Trends}\label{sec:aggregate}

To place the procurement premium in broader context, aggregate markup trends are constructed following De Loecker, Eeckhout, and Unger \cite{DeLoekerEeckhout2020}. For each year $t$, the sales-weighted aggregate markup is:
\begin{equation}\label{eq:dleu_agg}
    M_t^{sw} = \sum_i \omega_{it} \mu_{it}^A, \quad \omega_{it} = \frac{\text{sales}_{it}}{\sum_j \text{sales}_{jt}}.
\end{equation}
Figure~\ref{fig:dleu_aggregate} plots four specifications. Unlike the secular rise documented by DLEU for US manufacturing (1.21 in 1980 to 1.61 in 2016), Czech construction exhibits a \emph{declining} aggregate markup, from 2.45 in 2007 to 2.08 in 2021. The sales-weighted aggregate consistently exceeds the simple mean by 0.20--0.40 markup points, indicating that larger firms carry systematically higher markups.

\input{../input/tables/dleu_table_markup_specs.tex}

\begin{figure}[htbp]\centering
\includegraphics[width=0.85\textwidth]{dleu_fig1_aggregate_markup.pdf}
\caption{Aggregate Markups --- Czech Construction (CZ-NACE F)}\label{fig:dleu_aggregate}
\begin{minipage}{0.85\textwidth}\footnotesize
\emph{Notes:} Sales-weighted aggregate markup $M_t^{sw}$ under four specifications: ACF production-function-estimated $\hat{\theta}_{st}$ (benchmark), calibrated $\theta = 0.85$, firm-level cost share $\alpha^V_{it} = \text{COGS}_{it}/(\text{COGS}_{it} + r_t K_{it})$, and industry-median cost share. Vertical dashed lines mark Act~55/2012 (single-bid ban) and Act~134/2016 (MEAT criteria).
\end{minipage}
\end{figure}

Following DLEU (eq.~9), changes in the aggregate markup decompose into within-firm, market-share reallocation, cross, and net-entry components using the Haltiwanger (1997) demeaning:
\begin{equation}\label{eq:decomp}
    \Delta M_t = \underbrace{\sum_i m_{i,t-1} \Delta\mu_{it}}_{\Delta\text{within}} + \underbrace{\sum_i \tilde{\mu}_{i,t-1} \Delta m_{it}}_{\Delta\text{market share}} + \underbrace{\sum_i \Delta\mu_{it} \Delta m_{it}}_{\Delta\text{cross}} + \underbrace{\sum_{i \in E_t} \tilde{\mu}_{it} m_{it} - \sum_{i \in X_t} \tilde{\mu}_{i,t-1} m_{i,t-1}}_{\text{net entry}}.
\end{equation}
Three findings emerge (Figure~\ref{fig:dleu_decomp}, Table~\ref{tab:dleu_sectoral}). The within-firm component dominates the aggregate decline. Net entry contributes large year-to-year swings but is approximately neutral cumulatively. Most variation is within-industry, with reallocation between NACE sub-industries playing a secondary role---matching the DLEU finding that aggregate markup dynamics are driven by within-sector forces.

\begin{figure}[htbp]\centering
\includegraphics[width=0.85\textwidth]{dleu_fig4_decomposition.pdf}
\caption{Decomposition of Markup Growth at the Firm Level}\label{fig:dleu_decomp}
\begin{minipage}{0.85\textwidth}\footnotesize
\emph{Notes:} Cumulative decomposition following DLEU (2020, eq.~9). Each counterfactual starts from the 2007 aggregate markup and adds only the indicated component. ``Within'' holds market shares fixed; ``Reallocation'' holds markups fixed (market share + cross terms); ``Net entry'' isolates the contribution of entering and exiting firms.
\end{minipage}
\end{figure}

\input{../input/tables/dleu_table1_sectoral_decomp.tex}

The upper tail drives the aggregate (Figure~\ref{fig:dleu_percentiles}): the sales-weighted P90 exceeds 3.0 throughout, while the median is stable around 1.6--1.8. Splitting by procurement status (Figure~\ref{fig:dleu_pp_split}) confirms that procurement firms maintain systematically higher markups, with the gap narrowing after 2012 consistent with transparency reforms.

\begin{figure}[htbp]\centering
\includegraphics[width=0.85\textwidth]{dleu_fig3_distribution.pdf}
\caption{Percentiles of the Markup Distribution (sales-weighted)}\label{fig:dleu_percentiles}
\begin{minipage}{0.85\textwidth}\footnotesize
\emph{Notes:} Firms ranked by markup; percentiles computed using cumulative sales shares as weights (DLEU Figure~III, Panel~B). Average is the sales-weighted aggregate $M_t^{sw}$.
\end{minipage}
\end{figure}

\begin{figure}[htbp]\centering
\includegraphics[width=0.85\textwidth]{dleu_fig_procurement_split.pdf}
\caption{Aggregate Markup by Procurement Status}\label{fig:dleu_pp_split}
\begin{minipage}{0.85\textwidth}\footnotesize
\emph{Notes:} Sales-weighted aggregate markup computed separately for procurement ($pp = 1$) and non-procurement ($pp = 0$) firms. Vertical lines mark the 2012 and 2016 transparency reforms.
\end{minipage}
\end{figure}

When the decomposition is computed separately for procurement and non-procurement firms, the within-firm markup increase is concentrated among procurement firms. Firms entering government contracting experience markup increases, but the expanding procurement sector draws in firms with initially lower markups, partially offsetting the within-firm effect in the aggregate.

\iffalse % BMY figure disabled --- depends on disabled pipeline script
\begin{figure}[htbp]\centering
\includegraphics[width=0.85\textwidth]{bmy_decomposition.pdf}
\caption{Markup Decomposition: Within-Firm vs.\ Reallocation vs.\ Entry}\label{fig:decomposition}
\begin{minipage}{0.85\textwidth}\footnotesize
\emph{Notes:} Olley-Pakes \cite{OlleyPakes1996} style decomposition of the sales-weighted average markup. ``Within'' captures changes in markups holding market shares fixed. ``Reallocation'' captures the covariance between market share changes and markup levels. ``Net entry'' captures the contribution of entering and exiting firms.
\end{minipage}
\end{figure}
\fi


%% ============================================================
%% 5. THE PROCUREMENT MARKUP PREMIUM
%% ============================================================
\section{The Procurement Markup Premium}\label{sec:causal}

The cross-sectional results document a strong association between procurement and markups but do not establish causality. More productive firms may self-select into government tenders, and government contracts may provide stable demand that enables higher prices.

\subsection{Cross-Section and Panel Fixed Effects}\label{sec:panel_fe}

The primary specification exploits within-firm variation in procurement status across a panel of 1,267 firms:
\begin{equation}\label{eq:panel_fe}
    \ln\mu_{it} = \delta_1 \, pp_{it} + \alpha_i + \gamma_t + \nu_{it},
\end{equation}
where $\alpha_i$ and $\gamma_t$ are firm and year fixed effects. The coefficient $\delta_1$ is identified from the 474 switcher firms that change procurement status during the sample period. The estimate is $\hat{\delta}_1 = 0.131$ (SE = 0.005), corresponding to a 14.0\% premium. This is stable across specifications: pooled OLS with year$\times$NACE fixed effects gives 15.0\%, firm FE alone gives 14.8\%, and adding capital and COGS controls gives 14.1\%.

Robustness to omitted variable bias is assessed using Oster \cite{Oster2019}: the bias-adjusted coefficient under equal selection ($\delta = 1$) and $R_{\max} = \min(1.3\tilde{R}, 1)$ yields an identified set of $[0.116, 0.140]$, with $\delta^* = -5.9$---unobservables would need to work in the \emph{opposite} direction to explain away the premium.

A potential SUTVA concern is that procurement firms crowd out private demand for untreated firms. Using the KLMS~\cite{KLMS2025} estimate that government projects displace 27\% of private market output and the 30\% procurement share in Czech construction, the implied control-group markup depression is bounded at approximately 1--2 percentage points---meaningful but insufficient to explain the full premium.

\subsection{Entry and Exit Dynamics}\label{sec:entry_exit}

Following the exporter markup analysis of De Loecker and Warzynski \cite{DLW2012}, the procurement premium is decomposed into components explained by productivity selection and residual pricing power.

Firms entering public procurement experience an 11.5\% markup increase (Table \ref{tab:dlw_eval}); firms exiting see an 11.2\% decrease. The near-symmetric reversal is the strongest causal evidence: if the premium were driven by pre-existing firm characteristics, exit effects would be zero. The symmetry implies that procurement \emph{causes} markup changes, not merely selects high-markup firms.

\input{../input/tables/dlw_treatment_eval.tex}

\subsection{Event Study}\label{sec:event_study}

Figure~\ref{fig:event_study} plots the markup trajectory around procurement entry ($\tau = -3$ to $+5$, with $\tau = -1$ as reference). Pre-entry coefficients are flat and insignificant ($\tau = -3$: $-0.010$; $\tau = -2$: $-0.008$), validating the parallel trends assumption. The markup jumps sharply at entry ($\tau = 0$: $+0.134$) and remains persistently elevated through $\tau = +5$ ($\approx 0.09$--$0.11$). The absence of pre-trends combined with persistent post-entry effects is the cleanest causal evidence in the paper.

\begin{figure}[htbp]\centering
\includegraphics[width=0.85\textwidth]{dlw_event_study.pdf}
\caption{Event Study: Markup Dynamics Around Procurement Entry}\label{fig:event_study}
\begin{minipage}{0.85\textwidth}\footnotesize
\emph{Notes:} Coefficients from regressing $\log \mu^A_{it}$ on relative-year dummies ($\tau = -3$ to $+5$, $\tau = -1$ omitted as reference) with controls ($k$, cogs) and year $\times$ NACE FE. Firm-clustered 95\% CIs. Multi-switcher firms excluded (DLW fn.\ 52). Sample: firms with at least one year of pre-procurement data.
\end{minipage}
\end{figure}

\subsection{Productivity Decomposition}\label{sec:productivity_decomp}

Including translog productivity $\hat{\omega}_{it}$ as a control reduces the premium from 13.5\% to 3.2\% ($t = 4.1$). Productivity thus explains 76\% of the procurement-markup association---closely matching DLW's finding that productivity explains approximately 70\% of the export markup premium. The residual 3.2\% represents pricing power beyond cost advantages: procurement firms charge above what their productivity alone would predict, consistent with less elastic government demand or political favoritism \cite{BaranekTitl2024}.

Interacting the procurement dummy with procurement revenue share yields a coefficient of 0.037 ($t = 4.3$): each additional unit of procurement intensity raises markups by 3.7\%. This dose-response relationship reinforces the causal interpretation.

\subsection{Sub-Industry Heterogeneity}\label{sec:heterogeneity}

Interacting procurement status with NACE 2-digit dummies reveals heterogeneous premiums: buildings (NACE 41) 16.9\%, civil engineering (NACE 42) 15.5\%, and specialized construction (NACE 43) 7.8\%. The 2:1 ratio between buildings and specialized construction mirrors concentration differences and suggests the premium partly reflects market structure. By NACE sub-industry, the firm FE premium is largest in buildings (17.4\%, NACE 41) and smallest in civil engineering (6.4\%, NACE 42).

\subsection{Czech Transparency Reforms}\label{sec:policy}

The declining premium---from approximately 30\% in 2006 to approximately 10\% in 2021---coincides with the institutional reforms summarized in Table~\ref{tab:reforms}. Full-panel firm fixed effects estimates show a clear structural break around 2012, with the premium falling from approximately 0.30 (2006--2009) to 0.25 (2010--2015) to 0.10 (2016--2021).

An important caveat is that the decline may partly reflect compositional changes: as more firms enter procurement (the procurement share rose from 28\% to 55\%), the marginal entrants may be less productive and less able to extract rents, mechanically reducing the average premium. The within-firm decomposition partially addresses this by showing that within-firm markups of procurement firms also declined.

\iffalse % Favoritism subsection disabled --- depends on favoritism_decomposition pipeline script
\subsection{Favoritism versus Productivity}\label{sec:favoritism}

Baran\'{e}k and Titl \cite{BaranekTitl2024} show that political connections raise Czech procurement prices by 6.2\%, with no corresponding gains in project quality, and that EU co-funding oversight fully mitigates this overpricing. Their identification exploits electoral variation: when a connected politician loses power, the firm's contracts immediately return to normal prices.

\input{../input/tables/favoritism_decomposition.tex}

Table~\ref{tab:favoritism} decomposes the markup premium. Firms with above-median single-bid shares exhibit a larger premium than those facing greater competition (Panel~A). The premium declines from pre-2012 to post-2012, consistent with the single-bidding ban reducing favoritism (Panel~B). A calibration against Baran\'{e}k and Titl~\cite{BaranekTitl2024}, who estimate 6.2\% contract-price overpricing from political connections, places an upper bound of approximately 6 percentage points on the favoritism channel (Panel~C). Oster \cite{Oster2019} bounds yield $\delta^* = -6.05$, indicating unobservables would need to work in the opposite direction to explain away the premium (Panel~D). Both policy reforms are associated with significant declines in the procurement premium (Panel~E).

Panel~F provides a direct test using engineering cost estimates from the procurement register. Procuring authorities produce ex-ante cost estimates for each tender; the ratio of realized price to this estimate ($\mathrm{Rel\_Price}_{jt}$) measures overpricing at the contract level.\footnote{Engineering estimates explain 91\% of price variation in the full Datlab sample (Baran\'{e}k and Titl~\cite{BaranekTitl2024}, Table~1). I aggregate contract-level Rel\_Price to firm-year means, trimming to $(0.2, 5)$ following their methodology.} On the subsample of 385 procurement firms with estimate data, the baseline premium is 0.152 (row~1). Adding Rel\_Price as a control reduces it to only 0.141 (row~2)---overpricing accounts for less than 8\% of the premium. The Rel\_Price coefficient itself is small (0.021, $p < 0.01$): firms that bid higher relative to estimates have marginally higher markups, but the effect is economically negligible. Most strikingly, among procurement firms, the within-firm correlation between markups and Rel\_Price is essentially zero ($-0.002$, $p = 0.87$; row~3). High-markup firms do not overbid---they produce cheaply. The markup premium reflects productivity selection, not rent extraction through overpricing. The tender data lack EU co-funding indicators, precluding a direct test of the oversight channel identified by Baran\'{e}k and Titl~\cite{BaranekTitl2024}.
\fi

\iffalse % Selection on Observables disabled --- depends on lalonde pipeline script
\subsection{Selection on Observables}\label{sec:unconfound}

Following Imbens and Xu \cite{ImbensXu2024}, I estimate the ATT under the unconfoundedness assumption: conditional on pre-treatment covariates, procurement status is independent of potential outcomes:
\begin{equation}\label{eq:cia}
    W_i \perp Y_i(0) \mid X_i, \qquad 0 < \Pr(W_i = 1 \mid X_i) < 1.
\end{equation}

I estimate the ATT using the eleven-estimator toolkit of Imbens and Xu \cite{ImbensXu2024} on the full pooled cross-section (N=7,666, 3,443 treated). Covariates: $k_{it}$, cogs$_{it}$, employment, foreign-ownership dummy, market share, and NACE 2-digit industry dummies. Estimators span the taxonomy of selection-on-observables methods: difference-in-means, regression adjustment, outcome modeling (OM-Reg, OM-GRF), nearest-neighbor matching, propensity score matching, inverse probability weighting, covariate balancing propensity score (CBPS), entropy balancing, double machine learning (DML), and AIPW-GRF.

Table \ref{tab:lalonde} reports the results. All eleven estimators yield ATT estimates in the 0.14--0.19 range, with a median of approximately 0.15. The regression-adjusted estimators (OLS, OM-Reg, DML) cluster near 0.14--0.16, while matching-based methods (PSM, matching) and weighting-based methods (IPW, CBPS) yield slightly higher estimates (0.15--0.17). This convergence across methods with different assumptions about the propensity-score and outcome models provides robust evidence that the procurement premium is not an artifact of a particular identification strategy within the unconfoundedness paradigm.

\input{../input/tables/lalonde_comparison.tex}
\fi

\iffalse % Panel Data Methods disabled --- depends on fect pipeline script
\subsection{Panel Data Methods}\label{sec:panel}

The unconfoundedness design addresses selection on observables but requires the strong assumption that all relevant confounders are captured by pre-treatment covariates. Panel data methods can address selection on time-invariant unobservables.

\textbf{Two-way fixed effects.} On the balanced panel (26 firms $\times$ 16 years $=$ 416 observations, after dropping 13 always-treated units), the standard TWFE estimator with firm and year fixed effects yields $\hat{\tau} = 0.159$ (SE $= 0.022$). However, with staggered adoption, the TWFE estimator may be biased under heterogeneous treatment effects \cite{GoodmanBacon2021}. The Goodman-Bacon decomposition reveals substantial weight on ``later vs.\ earlier treated'' comparisons, motivating heterogeneity-robust alternatives.

\textbf{Heterogeneity-robust DiD.} The Borusyak, Jaravel, and Spiess \cite{BorusyakJaravel2024} imputation estimator yields an overall ATT of 0.142 (SE $= 0.013$), with the contemporaneous effect (period $t=0$) of 0.261 (SE $= 0.007$)---a large on-impact effect that fades to approximately 0.10 after several years.

\textbf{Synthetic difference-in-differences.} Cohort-specific SDID estimates confirm the declining premium pattern:

\begin{table}[htbp]\centering
\caption{SDID Cohort-Specific ATT Estimates}\label{tab:sdid}
\begin{threeparttable}
\begin{tabular}{lcc}
\toprule
Cohort (Entry Year) & ATT & SE \\
\midrule
2011 & 0.118 & 0.022 \\
2013 & 0.090 & 0.019 \\
2014 & 0.082 & 0.018 \\
2015 & 0.050 & 0.016 \\
2016 & 0.128 & 0.020 \\
2017 & $-$0.012 & 0.015 \\
2018 & 0.042 & 0.014 \\
\bottomrule
\end{tabular}
\begin{tablenotes}\footnotesize
\item \emph{Notes:} Synthetic DiD following Arkhangelsky et al.\ (2021). Balanced panel, 26 firms. Each row: staggered entry cohort with data-driven SC weights.
\end{tablenotes}
\end{threeparttable}
\end{table}

\textbf{Counterfactual estimators.} On the full unbalanced panel (N=7,666, 1,298 firms), I apply three counterfactual estimators from the \texttt{fect} R package \cite{LiuWangXu2022}: FEct (imputation TWFE), interactive fixed effects (IFEct, $r=2$ factors), and matrix completion (MC). All three yield ATT estimates in a tight range: FEct $= 0.125$ (SE $= 0.005$), IFEct $= 0.119$ (SE $= 0.011$), MC $= 0.122$ (SE $= 0.007$). These estimates are slightly below the firm FE estimate of 0.138 but within 1--2 standard errors, confirming the procurement premium is robust to latent-factor models of unobserved heterogeneity.

\input{../input/tables/fect_comparison.tex}

Dynamic treatment effects are flat and centered near zero during the pre-treatment period (Figure \ref{fig:counterfactual}). Placebo tests fail to reject the null of no pre-treatment effect, and no carryover effects are detected in the years following treatment exit.

\textbf{Sensitivity to parallel trend violations.} Following Rambachan and Roth \cite{RambachanRoth2023}, I compute robust confidence sets for the ATT under relative magnitude (RM) restrictions, which limit post-treatment parallel trend violations to at most $\bar{M}$ times the maximum pre-treatment violation. The breakdown value is $\bar{M} \approx 3.0$, indicating that the procurement premium is robust to PTA violations three times as large as those observed in pre-treatment periods.

\textbf{Augmented synthetic control.} The Ben-Michael, Feller, and Rothstein \cite{BenMichael2021} augmented SC method yields cohort-aggregated on-impact ATTs that decline over time: the 2011 and 2014 entry cohorts show effects of 0.26--0.30, while the 2017 and 2018 cohorts show effects of 0.04--0.08.

\begin{figure}[htbp]\centering
\includegraphics[width=0.85\textwidth]{s2_dyn.pdf}
\caption{Dynamic Treatment Effects: Counterfactual Estimators}\label{fig:counterfactual}
\begin{minipage}{0.85\textwidth}\footnotesize
\emph{Notes:} FEct, IFEct, and MC estimators from Liu, Wang, and Xu (2022). Panel (a): dynamic treatment effects with flat pre-trends. Panel (b): placebo tests (TOST rejects). Panel (c): carryover test (no persistence after exit). Full unbalanced panel, firms with $\geq$5 pre-treatment periods. Cluster-bootstrapped confidence intervals.
\end{minipage}
\end{figure}

Taken together, the four identification strategies yield a consistent picture: entering public procurement raises markups by approximately 14--15\%, with a large on-impact effect that declines over time. The full-panel firm FE estimate (14.0\%), the DLW entry effect (11.5\%), the design-based estimates (12--16\%), the counterfactual estimators (15.0\%), and the SDID cohort estimates (4--13\%) all point in the same direction. The convergence across methods with very different identifying assumptions---within-firm variation, selection on observables, synthetic control weighting, and matrix completion---provides strong evidence that the premium is not an artifact of a particular identification strategy.
\fi


%% ============================================================
%% 6. ROBUSTNESS
%% ============================================================
\section{Robustness}\label{sec:robustness}

\subsection{CWDL Extensions}\label{sec:cwdl}

Following Collard-Wexler and De Loecker \cite{CWDL2015}, six specifications vary the Markov process, survival correction, and instrument set. The procurement premium is 0.138 (SE = 0.005) under the baseline and all variants that include $pp_{it-1}$ in the Markov; it collapses to 0.003 (SE = 0.002) when procurement is excluded from the law of motion. The survival correction has negligible effects (Appendix Table~\ref{tab:cwdl}). Translog and OLS premiums fall within one standard error of the baseline.

\subsection{ADL Imperfect Competition}\label{sec:adl}

Under imperfect competition, input demand depends on the competitive environment, contaminating the scalar unobservable inversion at the heart of proxy variable estimators. Ackerberg and De Loecker \cite{ADL2024} show that this requires conditioning on a sufficient statistic---total competitor output under Cournot, an inclusive value under differentiated products.

The correction is implemented by defining markets as year $\times$ NACE 2-digit cells, including competitors' total deflated revenue in the first-stage polynomial and Markov transition, and adding oligopoly instruments to the moment conditions. The result is a null: the procurement premium changes by less than 0.005. Production function coefficients shift by 0.01--0.03, but the ratio determining the markup premium is approximately invariant.

Two explanations support this null. First, when competition is approximately symmetric within markets---firms face similar demand and the same competitor structure---the sufficient statistic is absorbed by market-level fixed effects, and the correction is irrelevant \cite{ADL2024}. In Czech construction, where firms routinely compete for both public and private projects within the same NACE$\times$year cell, this symmetry condition is plausible. Second, the KLMS \cite{KLMS2025} attenuation mechanism predicts that corrections to product and labor market power partially cancel in the markup differential across treated and control firms.

Table~\ref{tab:adl_instruments} reports a systematic instrument comparison. The output elasticity is stable at 1.01 across all accounting-based instruments but jumps to 1.5--1.9 with tender-level instruments, reflecting weak first-stage identification from bid counts. The premium is invariant at 0.170 across all 15 specifications---confirming identification from cost-share variation.

\input{../input/tables/adl_instrument_comparison.tex}

Caveats remain. The market definition is coarse; the Hansen $J$ test rejects across all overidentified specifications ($p < 0.001$), suggesting misspecification at this aggregation level.

\iffalse % Strong Exclusion disabled --- depends on ABGRS pipeline script
\subsection{Strong Exclusion Under Misspecification}\label{sec:abgrs}

Andrews, Barahona, Gentzkow, Rambachan, and Shapiro \cite{ABGRS2025} formalize when structural estimates are reliable under model misspecification. Their central result: estimators are ``approximately causally consistent'' only if they satisfy \emph{strong exclusion}---the moment conditions must rely on instruments that are mean-independent of the included controls. If this fails, misspecifying the control function corrupts the causal estimate even when the instruments are valid in the unconditional sense. Their Section~IV.D extends the framework directly to dynamic production function models with input endogeneity, the exact setting of ACF markup estimation.

I implement the ABGRS diagnostic in two steps. First, Table \ref{tab:strong_exclusion} reports the partial $R^2$ from regressing each ACF instrument on the control set (year $\times$ NACE 2-digit fixed effects). All instruments have $R^2 < 0.10$, with the largest value (0.098) for lagged procurement status, driven primarily by year effects (procurement expanded over the sample). Capital-based instruments load primarily on NACE effects ($R^2 \approx 0.04$), while COGS instruments are most orthogonal ($R^2 \approx 0.03$). These values indicate that ACF's instruments are largely mean-independent of the control set---strong exclusion approximately holds.

\input{../input/tables/strong_exclusion_scores.tex}

Second, I re-estimate the ACF following the ABGRS residualization recipe: I residualize the inputs ($k$, cogs, pp\_dummy) with respect to year $\times$ NACE fixed effects \emph{before} the estimator constructs its lagged-variable instruments, ensuring strong exclusion by construction. Table~\ref{tab:abgrs} reports the comparison. The residualized procurement premium is 0.140 (SE = 0.003), versus the baseline of 0.140 (SE = 0.003)---a difference of 0.03 log points, well within the baseline standard error. This confirms that the 14\% procurement markup premium is robust to the class of misspecifications ABGRS target: control-structure errors cannot corrupt the estimate because the instruments do not ride on control-variable variation.

\input{../input/tables/abgrs_residualization.tex}
\fi

\iffalse % Randomization Inference disabled --- depends on BH RI + Kitagawa pipeline scripts
\subsection{Randomization Inference}\label{sec:bh_ri}

Borusyak and Hull \cite{BorusyakHull2021, BorusyakHull2026} develop a framework for inference when treatment is a formula combining exogenous shocks with individual characteristics; see Borusyak, Hull, and Jaravel \cite{BorusyakHullJaravel2025} for a unified review of design-based formula instruments. Their approach: compare the observed test statistic to its distribution under counterfactual shock realizations. Applied to my setting, the counterfactual is the set of alternative assignments of procurement status across firms, holding fixed the year-level and industry-level procurement intensities.

Table~\ref{tab:bh_ri} reports randomization-inference confidence intervals for the procurement premium, constructed by permuting $pp_{it}$ either within each year (coarse) or within each (year, NACE 2-digit) stratum (fine), with $K = 1000$ permutations. The RI 95\% confidence intervals are $[0.134, 0.143]$ (coarse) and $[0.133, 0.144]$ (fine), essentially identical to the clustered-SE interval. The permuted premium distribution is centered at zero (SD $\approx 0.0025$) by construction, and none of the 1000 permutations in either strategy yield a premium exceeding the observed estimate of 0.138---corresponding to $p < 0.001$. This confirms that the 14\% premium cannot be attributed to the specific realization of contract assignments across firms within the observed strata.

\input{../input/tables/borusyak_hull_ri.tex}

\textbf{IV validity test.} Building on the BH recentering framework, I construct a shift-share instrument for procurement participation: $Z_{it} = \overline{pp}_i^{\text{pre}} \times \mathbf{1}[t \geq 2012]$, where $\overline{pp}_i^{\text{pre}}$ is firm $i$'s pre-reform procurement intensity and the shift is the 2012 single-bid ban. Recentering subtracts the expected instrument under counterfactual reform timings (2008--2016), yielding $\tilde{z}_{it} = Z_{it} - E[Z_{it} | X_i]$. The binarized instrument $\mathbf{1}[\tilde{z}_{it} > 0]$ has a first-stage $F = 287$ with $\Pr(D=1|Z=1) = 0.77$ vs.\ $\Pr(D=1|Z=0) = 0.41$.

I apply the Kitagawa \cite{Kitagawa2015} specification test, which tests whether the LATE IV validity conditions (exclusion, independence, monotonicity) are jointly consistent with the data by checking nonnegativity of the compliers' potential outcome densities. Table~\ref{tab:kitagawa} reports the variance-weighted KS statistic and bootstrap $p$-values for both the 2012 and 2016 reform instruments across three trimming constants $\xi \in \{0.01, 0.05, 0.10\}$. All $p$-values exceed 0.97---the test cannot reject IV validity for either reform instrument at any conventional significance level.

\input{../input/tables/kitagawa_iv_test.tex}
\fi

\iffalse % AGS/Andrews CS disabled --- depends on AGS, misspec, spec curve, and Armstrong pipeline scripts
\subsection{Transparency: AGS Sensitivity and Identification-Robust CS}\label{sec:ags_twostep}

Andrews, Gentzkow, and Shapiro \cite{AGS2017} and Andrews \cite{Andrews2025Lecture} argue that structural estimates should be reported alongside a sensitivity vector $\Lambda = -(G'WG)^{-1}G'W$ that maps moment-condition violations into parameter bias. I implement this diagnostic for ACF spec A on NACE 41 (the largest sub-industry, $N = 4{,}911$). Table~\ref{tab:ags_twostep_sensitivity} reports $\Lambda_{k,m}$ and $\Lambda_{\text{cogs},m}$ for each moment $m$. The lagged COGS moment is the most load-bearing for both production-function parameters, as expected: the ACF identification relies primarily on orthogonality of lagged input quantities with the productivity innovation.

\input{../input/tables/ags_twostep_sensitivity.tex}

I also report identification-robust confidence sets following Andrews \cite{Andrews2017ECMA}. The Wald CS assumes standard asymptotic normality; the S-statistic CS (Anderson-Rubin style) inverts the GMM criterion $N \cdot g(\beta)' W g(\beta)$ against the $\chi^2_{p,0.95}$ threshold without relying on the Wald approximation. Table~\ref{tab:ags_twostep_ci} shows that the S-statistic CS is 2.6--5 times wider than the Wald CS for the ACF production-function parameters, indicating moderate identification weakness in the first-stage GMM. This does not undermine the reduced-form procurement premium---the premium regression is an OLS on generated markups with a $t$-statistic of 32---but it does mean the ACF production-function coefficients themselves should be interpreted with wider uncertainty bounds than the analytical standard errors suggest.

\input{../input/tables/ags_twostep_ci.tex}

Andrews, Chen, and Tecchio \cite{AndrewsChenTecchio2026} show that under local misspecification, the Hansen $J$-statistic bounds the scope for ``weight hacking'': $\sqrt{J}$ measures how many standard errors a researcher can move the GMM estimate by varying the weighting matrix. All three ACF specifications in this paper are just-identified ($J = 0$, $df = 0$), so the weight matrix cannot move the estimate at all---a feature of the parsimonious Cobb-Douglas specification. Table~\ref{tab:misspec} further reports misspecification-robust standard errors for the procurement premium, computed via non-recentered cluster bootstrap (500 replications) following their recommendation. The bootstrap SE (0.005) matches the analytical clustered SE (0.005) almost exactly (ratio $\approx$ 1.01), confirming that misspecification does not inflate uncertainty in the reduced-form premium.

\input{../input/tables/misspecification_diagnostics.tex}

Figure~\ref{fig:spec_curve} aggregates all point estimates into a specification curve \cite{SimonsohnSimmonsNelson2020} covering ACF, OLS, cost-share, lalonde selection-on-observables (11 estimators), fect counterfactual methods, ABGRS residualization, Borusyak-Hull randomization inference, and Oster bounds. Out of 33 specifications, the procurement premium ranges from 0.003 (ACF without $pp_{it-1}$ in Markov) to 0.186 (raw difference-in-means), with a median of approximately 0.138. The preferred specification (ACF Base) falls near the center of the distribution, and the vast majority of specifications place the premium between 0.12 and 0.17.

\begin{figure}[htbp]\centering
\includegraphics[width=\textwidth]{specification_curve.pdf}
\caption{Specification Curve: Procurement Premium Across All Specifications}\label{fig:spec_curve}
\begin{minipage}{\textwidth}\footnotesize
\emph{Notes:} Each dot is the estimated procurement premium from one of 33 specifications spanning ACF/OLS production function estimates, selection-on-observables (lalonde), counterfactual estimators (fect), strong-exclusion residualization (ABGRS), randomization inference (Borusyak-Hull), and sensitivity bounds (Oster). Error bars are 95\% confidence intervals. The red dashed line marks the preferred estimate ($\hat{\beta}_{pp} = 0.138$).
\end{minipage}
\end{figure}

Armstrong \cite{Armstrong2025} formalizes this specification sensitivity through the \emph{expanded model approach}: rather than assuming the baseline model is correct, expand to a larger class of specifications and report the worst-case bias within this class. Table~\ref{tab:armstrong_ci} implements this diagnostic. The standard firm-clustered 95\% CI is $[0.132, 0.144]$. The Armstrong bias-adjusted CI---accounting for the maximum specification-curve deviation as an upper bound on bias---widens to $[0.052, 0.224]$ when restricting to internally consistent specifications (those including $pp_{it-1}$ in the Markov per De Loecker 2013), or $[-0.003, 0.279]$ over the full expanded model. Even the widest interval excludes zero only marginally, but the interquartile range of the specification curve $[0.132, 0.150]$ indicates that the central tendency is tightly concentrated around 14\%.

\input{../input/tables/armstrong_bias_adjusted_ci.tex}
\fi

\iffalse % Cross-Industry Validation disabled --- depends on spec curve pipeline script
\subsection{Cross-Industry Validation}\label{sec:cross_industry_rob}

A natural concern is whether the procurement premium is specific to Czech construction. Using the Orbis all-industry panel (Section \ref{sec:data}), I estimate ACF production functions and markup premiums for the five industries with the highest procurement participation. The construction premium (13.5\%, pooled NACE 41--43) matches the MagnusWeb baseline (14.0\%) on an independent dataset with $16\times$ more firms (24,382 vs.\ 1,267). Knowledge-intensive services exhibit even larger premiums: architecture and engineering (16.6\%, NACE 71) and IT services (15.3\%, NACE 62). Wholesale trade, where products are more homogeneous, shows a smaller premium (4.8\%, NACE 46). The cross-industry pattern is consistent with procurement conferring larger markup advantages in industries with differentiated products and high fixed costs.
\fi

\subsection{First-Stage and Markov Sensitivity}\label{sec:firststage}

Polynomial order in the first stage (2--5) has negligible effects on the premium: all estimates fall within 0.136--0.140. Enriching the Markov process with higher-order terms in $\omega_{it-1}$ or interactions with NACE dummies adds less than 0.5 percentage points to the Markov $R^2$ and leaves the premium invariant. DGM cross-validation against the ProdFun implementation confirms that implementation differences are negligible relative to specification sensitivity (Appendix~\ref{app:dgm}).


%% ============================================================
%% 7. CONCLUSION
%% ============================================================
\section{Conclusion}\label{sec:conclusion}

This paper establishes three findings on markups and public procurement in Czech construction. First, the procurement markup premium is approximately 14\%, robust across nine estimation methods, functional forms, and multiple identification strategies; an event study shows flat pre-trends and a sharp entry effect, and productivity explains 76\% of the markup association. Second, the premium declined from 30\% in 2006 to 10\% in 2021, with Act 55/2012 (single-bid ban) and Act 134/2016 (MEAT criteria) associated with DiD estimates of $-2.6$ and $-2.2$ log points. Third, aggregate sales-weighted markups in Czech construction, unlike the secular rise documented by De Loecker, Eeckhout, and Unger \cite{DeLoekerEeckhout2020} for US manufacturing, exhibit a declining trajectory driven primarily by within-firm changes.

Two limitations qualify these results. The sensitivity of markup estimates to the variable input definition---COGS versus intermediate inputs---remains unresolved empirically. MagnusWeb employment brackets prevent estimation of the labor supply elasticity needed for full KLMS-style double market power analysis; continuous employment data from Orbis or CZSO would address this. Bid-level data from the Czech procurement register---currently limited to winning bids---would enable structural rent and incidence calculations that could distinguish market power from cost efficiency more sharply.

The policy implication is direct: Czech transparency reforms reduced the procurement markup premium by roughly two-thirds over 15 years. The 2012 single-bidding ban and the 2016 MEAT criteria shift are both associated with significant declines. These reforms provide a template for other countries with low procurement competition.


%% ============================================================
%% APPENDIX A
%% ============================================================
\appendix
\section{Methods}\label{app:methods}

\subsection{Markup Formula Derivation}\label{app:markup_derivation}

Consider a firm maximizing profits by choosing variable input $V_{it}$:
\begin{equation}
    \max_{V_{it}} \quad P_{it} Q_{it} - P^V_{it} V_{it} - C(K_{it}),
\end{equation}
where $Q_{it} = F(V_{it}, K_{it}) \exp(\omega_{it})$ is the production function. The first-order condition with respect to $V_{it}$ is:
\begin{equation}
    P_{it} \frac{\partial Q_{it}}{\partial V_{it}} = \mu_{it}^{-1} \cdot P_{it} \cdot \frac{\partial Q_{it}}{\partial V_{it}} \cdot \mu_{it} = P^V_{it},
\end{equation}
which can be rearranged to:
\begin{equation}
    \mu_{it} \equiv \frac{P_{it}}{\lambda_{it}} = \frac{\theta^V_{it}}{\alpha^V_{it}},
\end{equation}
where $\theta^V_{it} = \frac{\partial \ln Q_{it}}{\partial \ln V_{it}}$ is the output elasticity of the variable input and $\alpha^V_{it} = \frac{P^V_{it} V_{it}}{P_{it} Q_{it}}$ is the expenditure share of the variable input in revenue. Under perfect competition $\mu_{it} = 1$ and $\theta^V_{it} = \alpha^V_{it}$; deviations from unity indicate market power.

\subsection{ACF Estimation}\label{app:acf}

The ACF estimator proceeds in two stages.

\textbf{Stage 1.} Estimate expected output as a flexible function of observables:
\begin{equation}
    y_{it} = \phi_t(v_{it}, k_{it}, pp_{it}) + \varepsilon_{it},
\end{equation}
where $\phi_t(\cdot)$ is approximated by a third-order polynomial. This yields $\hat{\phi}_{it} = y_{it} - \hat{\varepsilon}_{it}$.

\textbf{Stage 2.} For a candidate parameter vector $\beta$, compute:
\begin{equation}
    \hat{\omega}_{it}(\beta) = \hat{\phi}_{it} - f(v_{it}, k_{it}; \beta),
\end{equation}
and estimate the Markov process:
\begin{equation}
    \hat{\omega}_{it} = \rho \hat{\omega}_{it-1} + \gamma pp_{it-1} + \delta \hat{p}_{it-1} + \xi_{it},
\end{equation}
where $\hat{p}_{it-1}$ is the predicted survival probability from a probit model. The innovation $\hat{\xi}_{it}(\beta)$ is orthogonal to the instruments $\mathbf{Z}_{it}$:
\begin{equation}
    \E[\hat{\xi}_{it}(\beta) \cdot \mathbf{Z}_{it}] = 0.
\end{equation}
The GMM criterion is minimized over $\beta$ using Nelder-Mead followed by BFGS refinement.

\textbf{Standard errors.} I compute analytical standard errors following Ackerberg, Chen, and Hahn \cite{ACH2012}. The GMM sandwich variance is:
\begin{equation}
    \hat{V} = \frac{1}{N} (G'WG)^{-1} G'W \hat{S} W G (G'WG)^{-1},
\end{equation}
where $G$ is the Jacobian of moment conditions, $W$ is the weighting matrix, and $\hat{S}$ is the firm-clustered meat matrix. Standard errors for markups are obtained via the delta method.

\subsection{Translog Output Elasticity}\label{app:translog}

Under the translog specification (Equation \ref{eq:tl}), the output elasticity of the variable input is firm- and time-specific:
\begin{equation}
    \theta^V_{it} = \hat{\beta}_v + 2\hat{\beta}_{vv} v_{it} + \hat{\beta}_{vk} k_{it}.
\end{equation}
The markup is then $\mu_{it} = \theta^V_{it} / \alpha^V_{it}$, where the expenditure share $\alpha^V_{it}$ is computed directly from the data. Under Cobb-Douglas, $\theta^V_{it} = \hat{\beta}_v$ is constant across firms, and markup variation comes entirely from the expenditure share.

\subsection{Moment Conditions}\label{app:moments}

The baseline instrument set for the Cobb-Douglas specification is $\mathbf{Z}_{it} = (1, k_{it}, \text{cogs}_{it-1})$, yielding a just-identified system with 3 moments and 3 parameters ($\beta_0, \beta_k, \beta_{\text{cogs}}$). For the translog, the instrument set expands to include squares and interactions: $\mathbf{Z}_{it} = (1, k_{it}, \text{cogs}_{it-1}, k_{it}^2, \text{cogs}_{it-1}^2, k_{it} \cdot \text{cogs}_{it-1})$, yielding a just-identified system with 6 moments and 6 parameters.

For overidentification tests, I add $k_{it-1}$ and $\text{cogs}_{it-2}$ to the instrument set, providing 2 additional moments. The Hansen $J$ statistic tests the overidentifying restrictions.

%% ============================================================
%% APPENDIX B
%% ============================================================
\section{Additional Results}\label{app:additional}

\subsection{Markup Summary Statistics by Sub-Industry}\label{app:markup_nace}

\begin{table}[htbp]\centering
\caption{Markup Distribution by NACE 2-Digit Sub-Industry}\label{tab:markup_nace}
\begin{threeparttable}
\begin{tabular}{lccccc}
\toprule
 & Mean & SD & P10 & P50 & P90 \\
\midrule
NACE 41 (Buildings) & 2.14 & 0.23 & 1.82 & 2.17 & 2.41 \\
NACE 42 (Civil Eng.) & 2.35 & 0.33 & 2.01 & 2.29 & 2.82 \\
NACE 43 (Specialized) & 1.71 & 0.26 & 1.47 & 1.65 & 2.04 \\
\midrule
All & 1.99 & 0.35 & 1.53 & 1.94 & 2.44 \\
\bottomrule
\end{tabular}
\begin{tablenotes}\footnotesize
\item \emph{Notes:} Baseline specification: ACF Cobb-Douglas with survival correction and procurement dummy in the Markov transition. $N = 7{,}666$ (observations with valid markups). Markups expressed in levels; log markups used in regressions.
\end{tablenotes}
\end{threeparttable}
\end{table}

\iffalse % Premium by Year disabled --- depends on premium_timeseries pipeline script
\subsection{Procurement Premium by Year}\label{app:premium_year}

The cross-sectional procurement premium (the coefficient $\hat{\beta}_1$ in Equation \ref{eq:cross} estimated year by year) declines monotonically from 0.30 in 2006 to 0.10 in 2021. The decline accelerates around 2012, coinciding with the single-bidding ban (Act 55/2012).
\fi

\subsection{CWDL Robustness Specifications}\label{app:cwdl}

Following Collard-Wexler and De Loecker \cite{CWDL2015}, I estimate six specifications varying the Markov process and the instrument set:

\begin{table}[htbp]\centering
\caption{CWDL Robustness: Procurement Premium Across Markov Specifications}\label{tab:cwdl}
\begin{threeparttable}
\begin{tabular}{lcccc}
\toprule
Specification & Raw Premium & SE & Reg.\ Premium & SE \\
\midrule
Base (survival + pp in Markov) & 0.186 & 0.003 & 0.138 & 0.005 \\
No survival correction & 0.187 & 0.003 & 0.138 & 0.005 \\
No pp in Markov & 0.075 & 0.003 & 0.003 & 0.002 \\
Plain ACF & 0.077 & 0.003 & 0.003 & 0.002 \\
Translog & 0.173 & 0.003 & 0.139 & 0.005 \\
OLS & 0.195 & 0.008 & 0.140 & 0.009 \\
\bottomrule
\end{tabular}
\begin{tablenotes}\footnotesize
\item \emph{Notes:} $N = 7{,}666$ (ACF) or $N = 9{,}164$ (OLS). Regression controls: $k$, cogs, year$\times$NACE FEs. SEs clustered by firm.
\end{tablenotes}
\end{threeparttable}
\end{table}

\subsection{ADL Imperfect Competition: Detailed Results}\label{app:adl_detail}

Under the ADL correction, the production function coefficients shift but the premium is invariant:

\begin{itemize}[nosep]
    \item NACE 41: $\hat{\beta}_{\text{cogs}}$ shifts from 0.957 (baseline) to 0.934 (ADL). Premium: 0.138 $\to$ 0.136.
    \item NACE 42: $\hat{\beta}_{\text{cogs}}$ shifts from 1.100 to 1.078. Premium: 0.140 $\to$ 0.137.
    \item NACE 43: $\hat{\beta}_{\text{cogs}}$ shifts from 0.927 to 0.910. Premium: 0.138 $\to$ 0.134.
\end{itemize}

The Hansen $J$ test rejects for NACE 41 ($p = 0.004$) and is marginal for NACE 43 ($p = 0.07$), suggesting that the coarse market definition may introduce misspecification. NACE 42 passes ($p = 0.87$).

\iffalse % Full Unconfoundedness disabled --- depends on lalonde pipeline script
\subsection{Full Unconfoundedness Results}\label{app:full_unconfound}

\begin{table}[htbp]\centering
\caption{Full Battery of Unconfoundedness ATT Estimates}\label{tab:full_unconfound}
\begin{threeparttable}
\begin{tabular}{lcccc}
\toprule
 & \multicolumn{2}{c}{Contract Year} & \multicolumn{2}{c}{Placebo} \\
\cmidrule(lr){2-3}\cmidrule(lr){4-5}
Method & ATT & SE & ATT & SE \\
\midrule
Difference-in-Means & 0.12 & 0.02 & 0.03 & 0.02 \\
OLS Regression & 0.16 & 0.01 & $-$0.00 & 0.01 \\
Ridge Regression & 0.15 & 0.01 & 0.01 & 0.01 \\
GRF & 0.13 & 0.01 & 0.03 & 0.01 \\
1-NN Matching & 0.14 & 0.02 & 0.01 & 0.02 \\
PSM & 0.13 & 0.01 & $-$0.00 & 0.01 \\
IPW & 0.15 & 0.01 & 0.00 & 0.01 \\
AIPW & 0.15 & 0.01 & 0.00 & 0.01 \\
AIPW-GRF & 0.15 & 0.01 & 0.00 & 0.01 \\
DML-ElasticNet & 0.16 & 0.01 & $-$0.01 & 0.01 \\
DML-RF & 0.14 & 0.01 & 0.01 & 0.01 \\
\bottomrule
\end{tabular}
\begin{tablenotes}\footnotesize
\item \emph{Notes:} Balanced panel of 26 firms, 312 observations. Treatment: first procurement contract year. Placebo: one year before actual contract. All methods use 4 lags of outcomes and covariates.
\end{tablenotes}
\end{threeparttable}
\end{table}
\fi

\iffalse % Panel Data Event Study disabled --- depends on fect pipeline script
\subsection{Panel Data Event Study}\label{app:event_study}

The Borusyak, Jaravel, and Spiess \cite{BorusyakJaravel2024} event study estimates show a clear pattern: pre-treatment coefficients are indistinguishable from zero (supporting parallel trends), the on-impact effect at $t=0$ is approximately 0.26, and the effect stabilizes around 0.10--0.14 after 3--5 years.
\fi

\iffalse % HonestDiD disabled --- depends on fect pipeline script
\subsection{Sensitivity to HonestDiD}\label{app:honestdid}

I assess sensitivity to violations of parallel trends using the Rambachan and Roth \cite{RambachanRoth2023} HonestDiD framework. Under the DeltaRM restriction (bounded relative magnitudes of pre-trend violations), the confidence set for the ATT remains strictly positive for $\bar{M} > 4$ in both the full window (9 pre, 11 post) and the short window (3 pre, 4 post), indicating that the result is robust to moderate violations of parallel trends.
\fi

\subsection{DGM Validation and Revenue Bias}\label{app:dgm}

I cross-validate all estimates against the De Ridder, Grassi, and Morzenti \cite{DGM2026} Python implementation (ProdFun class), which provides an independent codebase for ACF estimation. The CD estimates match to within 0.001 for the COGS elasticity and 0.002 for the markup mean across all three NACE 2-digit industries. This confirms that implementation differences---optimizer choice, polynomial order, convergence criteria---are negligible relative to specification sensitivity.

A broader concern is that revenue-based markup estimates may be biased because financial data record revenue, not physical output \cite{Bond2021,KletteGriliches1996}. DGM \cite{DGM2026} address this directly. Under Cobb-Douglas with common technology within sectors, the revenue-to-quantity bias is a sector-level constant that shifts markup \emph{levels} but preserves all cross-firm \emph{variation}, including dispersion and trends. Their simulations confirm a 0.94 correlation between revenue-based and true markups, and French production data validate this at 0.80 in first differences. Since my analysis focuses on the procurement markup \emph{premium}---a cross-firm difference---rather than the average markup level, the revenue proxy does not compromise the treatment effect.

DGM also show that if the variable input faces a wedge---such as adjustment costs or monopsony power---the production approach measures the product of the markup and the input wedge (their Section 2.4.1). This provides an additional justification for using COGS rather than labor as the variable input: labor markets may feature monopsony power \cite{KLMS2025}, which would contaminate labor-based markups with a wage markdown. Materials markets, by contrast, are plausibly competitive in construction, making COGS the cleaner measure of the output elasticity relevant for markup estimation.

\subsection{Sample Restriction Sensitivity}\label{app:sample_robust}

The rebuilt sample (1,521 firms, 9,164 observations) differs from the original thesis sample (1,297 firms, 7,261 observations) due to corrected merges. To verify that sample construction does not drive results, I compare estimates on the overlap sample (9,087 observations in both datasets) and the 77 rebuilt-only observations.

The rebuilt-only observations have slightly lower mean log sales (17.74 vs.\ 18.13) and a substantially lower procurement share (22.1\% vs.\ 42.4\%). The production function estimates and the procurement premium are virtually unchanged when excluding these observations, confirming that the sample correction improves data quality without affecting the substantive conclusions.

\iffalse % BMY figure disabled --- depends on disabled pipeline script
\begin{figure}[htbp]\centering
\includegraphics[width=0.85\textwidth]{bmy_sample_sensitivity.pdf}
\caption{Sample Restriction Sensitivity: Original vs.\ Rebuilt}\label{fig:sample_sensitivity}
\begin{minipage}{0.85\textwidth}\footnotesize
\emph{Notes:} Comparison of markup distributions and procurement premiums across the original (1,297 firms) and rebuilt (1,521 firms) samples. The 77 rebuilt-only observations are smaller firms with lower procurement exposure.
\end{minipage}
\end{figure}
\fi

\iffalse % Raval Overidentification disabled --- depends on raval_test pipeline script
\subsection{Raval Overidentification}\label{app:raval}

Under cost minimization, markups estimated from \emph{any} flexible input should coincide: $\mu_{it} = \theta^X_{it}/\alpha^X_{it}$ for all variable inputs $X$. Raval \cite{Raval2023} tests this overidentifying restriction across seven datasets and finds that labor-based and materials-based markups are \emph{negatively} correlated everywhere, rejecting Hicks-neutral productivity. I replicate this test for Czech construction using three inputs: labor (wages), materials (COGS), and a composite (wages + COGS). Regressing log labor markups on log COGS markups with industry-year fixed effects, I find $\hat{\beta} = 0.28$--$0.82$---positive and substantially different from Raval's negative correlations. While $\beta = 1$ is rejected, the positive sign means the two markup measures \emph{agree on firm rankings}, which is the relevant property for treatment evaluation. De Ridder, Grassi, and Morzenti \cite{DGM2026} show formally (their Section 2.4.1) that when a variable input faces a wedge (e.g., monopsony), the production approach measures the product of the markup and the input wedge. The positive Czech $\beta$ suggests the wedge is monotone rather than contradictory, so the treatment effect is preserved even when overidentification fails.
\fi

\subsection{GNR Non-Identification}\label{app:gnr}

Gandhi, Navarro, and Rivers \cite{GNR2020} prove that proxy variable methods---including ACF---generically fail to identify the flexible input elasticity for gross output production functions. The flexible input appears in both the production function and the control function, creating collinearity that prevents separate identification. Their solution is a share regression: under perfect competition, the revenue share of the flexible input equals its output elasticity.

In my data, the OLS coefficient on the variable input is virtually identical to the ACF coefficient: the maximum discrepancy across the three NACE 2-digit industries is 3.4\% (NACE 43: OLS $\hat{\beta}_{\text{cogs}} = 0.899$ vs.\ ACF $\hat{\beta}_{\text{cogs}} = 0.927$). This near-equivalence is exactly what the GNR theorem predicts. Importantly, while the \emph{level} of the COGS elasticity may be unidentified, the \emph{procurement premium ratio}---which depends on the ratio $\theta^V_{it}/\alpha^V_{it}$ rather than the level of $\theta^V_{it}$---survives even if the production function coefficient is biased, provided the bias is common to treated and control firms. The specification sensitivity analysis (Section \ref{sec:sensitivity}) confirms this: the regression-adjusted premium is 0.140 under both OLS and ACF.

\iffalse % Double Market Power appendix disabled --- depends on Orbis data
\subsection{Double Market Power: Brief Benchmark}\label{app:klms}

Following the Kroft, Luo, Mogstad, and Setzler \cite{KLMS2025} framework, the revenue elasticity $(1-\epsilon)$ is estimated from the regression of log revenues on log intermediate expenditures among non-procurement firms with firm fixed effects: $(1-\epsilon) = 0.90$ for Czech construction (price markup $1/(1-\epsilon) = 1.12$), comparable to the KLMS US estimate of $0.86$ (price markup $1.16$). The inverse labor supply elasticity $\hat{\theta}$ requires continuous employment data not available in MagnusWeb; with a calibrated $\theta = 0.13$ (matching the KLMS US construction estimate), the implied double wage markdown is $(1+\theta)^{-1} \times (1-\epsilon) = 0.89 \times 0.90 = 0.80$. Full identification of $\theta$ from Czech data is blocked on Orbis or CZSO continuous employment access.
\fi

%% ============================================================
%% BIBLIOGRAPHY
%% ============================================================
\newpage
\begingroup
\small
\begin{thebibliography}{99}

\bibitem{ABGRS2025} Andrews, I., N.~Barahona, M.~Gentzkow, A.~Rambachan, and J.M.~Shapiro (2025). ``Structural Estimation Under Misspecification: Theory and Implications for Practice,'' \emph{Quarterly Journal of Economics}, forthcoming.

\bibitem{AGS2017} Andrews, I., M.~Gentzkow, and J.M.~Shapiro (2017). ``Measuring the Sensitivity of Parameter Estimates to Estimation Moments,'' \emph{Quarterly Journal of Economics}, 132(4), 1553--1592.

\bibitem{Andrews2017ECMA} Andrews, I. (2017). ``Valid Two-Step Identification-Robust Confidence Sets for GMM,'' \emph{Econometrica}, 84(6), 2155--2182.

\bibitem{Andrews2025Lecture} Andrews, I. (2025). ``Recent Developments in Econometrics for Transparency and Communication,'' AEA Continuing Education Lecture.

\bibitem{Armstrong2025} Armstrong, T.B. (2025). ``Misspecification in Econometrics: A Selective Review,'' Working Paper, USC.

\bibitem{AndrewsChenTecchio2026} Andrews, I., J.~Chen, and O.~Tecchio (2026). ``The Purpose of an Estimator Is What It Does: Misspecification, Estimands, and Over-Identification,'' Working Paper, arXiv:2508.13076v5.

\bibitem{SimonsohnSimmonsNelson2020} Simonsohn, U., J.P.\ Simmons, and L.D.\ Nelson (2020). ``Specification Curve Analysis,'' \emph{Nature Human Behaviour}, 4, 1208--1214.

\bibitem{ACF2015} Ackerberg, D., K.~Caves, and G.~Frazer (2015). ``Identification Properties of Recent Production Function Estimators,'' \emph{Econometrica}, 83(6), 2411--2451.

\bibitem{ACH2012} Ackerberg, D., X.~Chen, and J.~Hahn (2012). ``A Practical Asymptotic Variance Estimator for Two-Step Semiparametric Estimators,'' \emph{Review of Economics and Statistics}, 94(4), 938--952.

\bibitem{ADL2024} Ackerberg, D.\ and J.~De Loecker (2024). ``Production Function Identification Under Imperfect Competition,'' CEPR Discussion Paper No.~19640.

\bibitem{AngristPischke2009} Angrist, J.D.\ and J.-S.~Pischke (2009). \emph{Mostly Harmless Econometrics: An Empiricist's Companion}, Princeton University Press.

\bibitem{ArkhangelskySynthDiD} Arkhangelsky, D., S.~Athey, D.A.~Hirshberg, G.W.~Imbens, and S.~Wager (2021). ``Synthetic Difference-in-Differences,'' \emph{American Economic Review}, 111(12), 4088--4118.

\bibitem{Athey2021} Athey, S., M.~Bayati, N.~Doudchenko, G.~Imbens, and K.~Khosravi (2021). ``Matrix Completion Methods for Causal Panel Data Models,'' \emph{Journal of the American Statistical Association}, 116(536), 1716--1730.

\bibitem{BaranekTitl2024} Baran\v{e}k, B.\ and V.~Titl (2024). ``The Cost of Favoritism in Public Procurement,'' \emph{Journal of Law and Economics}, 67(2), 445--477.

\bibitem{BenMichael2021} Ben-Michael, E., A.~Feller, and J.~Rothstein (2021). ``The Augmented Synthetic Control Method,'' \emph{Journal of the American Statistical Association}, 116(536), 1789--1803.

\bibitem{BMY2026} Benkard, C.L., A.~Miller, and A.~Yurukoglu (2026). ``The Rise of Market Power: Comment,'' Working Paper.

\bibitem{Bond2021} Bond, S., A.~Hashemi, G.~Kaplan, and P.~Zoch (2021). ``Some Unpleasant Markup Arithmetic,'' \emph{Journal of Monetary Economics}, 121, 1--14.

\bibitem{BorusyakHull2021} Borusyak, K.\ and P.~Hull (2021). ``Non-Random Exposure to Exogenous Shocks: Theory and Applications,'' NBER Working Paper 29513.

\bibitem{BorusyakHull2026} Borusyak, K.\ and P.~Hull (2026). ``Optimal Formula Instruments,'' Working Paper.

\bibitem{BorusyakHullJaravel2025} Borusyak, K., P.~Hull, and X.~Jaravel (2025). ``Design-Based Identification with Formula Instruments: A Review,'' \emph{Econometrics Journal}, 28, 83--108.

\bibitem{BorusyakJaravel2024} Borusyak, K., X.~Jaravel, and J.~Spiess (2024). ``Revisiting Event-Study Designs: Robust and Efficient Estimation,'' \emph{Review of Economic Studies}, 91(6), 3253--3285.

\bibitem{CinelliHazlett2020} Cinelli, C.\ and C.~Hazlett (2020). ``Making Sense of Sensitivity: Extending Omitted Variable Bias,'' \emph{Journal of the Royal Statistical Society: Series B}, 82(1), 39--67.

\bibitem{CWDL2015} Collard-Wexler, A.\ and J.~De Loecker (2015). ``Reallocation and Technology: Evidence from the US Steel Industry,'' \emph{American Economic Review}, 105(1), 131--171.

\bibitem{CWDLKvsOmega} Collard-Wexler, A.\ and J.\ De Loecker (2020). ``Productivity and Capital Measurement Error,'' \emph{NBER Working Paper} 22437.

\bibitem{CMOY2023} Conlon, C., N.H.~Miller, T.~Otgon, and Y.~Yao (2023). ``Rising Markups, Rising Prices?'' \emph{AEA Papers and Proceedings}, 113, 279--283.

\bibitem{DeLoekerEeckhout2020} De Loecker, J., J.~Eeckhout, and G.~Unger (2020). ``The Rise of Market Power and the Macroeconomic Implications,'' \emph{Quarterly Journal of Economics}, 135(2), 561--644.

\bibitem{DLEU2025reply} De Loecker, J., J.~Eeckhout, and G.~Unger (2025). ``The Rise of Market Power and the Macroeconomic Implications: Reply to Benkard, Miller, and Yurukoglu,'' NBER Working Paper.

\bibitem{DeLoecker2013} De Loecker, J. (2013). ``Detecting Learning by Exporting,'' \emph{American Economic Journal: Microeconomics}, 5(3), 1--21.

\bibitem{DLW2012} De Loecker, J.\ and F.~Warzynski (2012). ``Markups and Firm-Level Export Status,'' \emph{American Economic Review}, 102(6), 2437--2471.

\bibitem{DeLoeckerScott2025} De Loecker, J. and P.T.\ Scott (2025). ``Markup Estimation using Production and Demand Data: An Application to the US Brewing Industry,'' Working Paper, NYU Stern / KU Leuven.

\bibitem{Demirer2025} Demirer, M.\ (2025). ``Production Function Estimation with Factor-Augmenting Technology: An Application to Markups,'' \emph{Econometrica}, conditionally accepted.

\bibitem{DiezLeighTambunlertchai2018} Diez, F., D.~Leigh, and S.~Tambunlertchai (2018). ``Global Market Power and its Macroeconomic Implications,'' \emph{IMF Working Paper} 18/137.

\bibitem{DoraszelskiJaumandreu2018} Doraszelski, U.\ and J.~Jaumandreu (2018). ``Measuring the Bias of Technological Change,'' \emph{Journal of Political Economy}, 126(3), 1027--1084.

\bibitem{FlynnHuoKako2019} Flynn, Z., V.~Huo, and J.~Kako (2019). ``Identifying Production Functions with Imperfect Competition,'' Working Paper.

\bibitem{DGM2026} De Ridder, M., B.~Grassi, and G.~Morzenti (2026). ``The Hitchhiker's Guide to Markup Estimation,'' \emph{Econometrica}, 94(1), 137--168.

\bibitem{GNR2020} Gandhi, A., S.~Navarro, and D.~Rivers (2020). ``On the Identification of Gross Output Production Functions,'' \emph{Journal of Political Economy}, 128(8), 2973--3016.

\bibitem{GoodmanBacon2021} Goodman-Bacon, A.\ (2021). ``Difference-in-Differences with Variation in Treatment Timing,'' \emph{Econometrica}, 89(5), 2261--2298.

\bibitem{Hall1988} Hall, R.E.\ (1988). ``The Relation Between Price and Marginal Cost in U.S.\ Industry,'' \emph{Journal of Political Economy}, 96(5), 921--947.

\bibitem{Hall2018} Hall, R.E.\ (2018). ``New Evidence on the Markup of Prices over Marginal Costs and the Role of Mega-Firms in the US Economy,'' \emph{Journal of Economic Perspectives}, 32(3), 53--72.

\bibitem{ImbensXu2024} Imbens, G.W.\ and Y.~Xu (2024). ``LaLonde (1986) After Nearly Four Decades: Lessons Learned,'' Working Paper.

\bibitem{PalgutaPertold2017} Palguta, J.\ and F.~Pertold-Gebicka (2017). ``Manipulation of Procurement Contracts: Evidence from the Introduction of Discretionary Thresholds,'' \emph{American Economic Journal: Economic Policy}, 9(2), 293--315.

\bibitem{RambachanRoth2023} Rambachan, A.\ and J.~Roth (2023). ``A More Credible Approach to Parallel Trends,'' \emph{Review of Economic Studies}, 90(5), 2555--2591.

\bibitem{Raval2023} Raval, D.\ (2023). ``Testing the Production Approach to Markup Estimation,'' \emph{Review of Economic Studies}, 90(5), 2592--2611.

\bibitem{KLMS2025} Kroft, K., Y.~Luo, M.~Mogstad, and B.~Setzler (2025). ``Imperfect Competition and Rents in Labor and Product Markets: The Case of the Construction Industry,'' \emph{American Economic Review}, 115(9), 2926--2969.

\bibitem{Kitagawa2015} Kitagawa, T.\ (2015). ``A Test for Instrument Validity,'' \emph{Econometrica}, 83(5), 2043--2063.

\bibitem{KletteGriliches1996} Klette, T.~J.\ and Z.~Griliches (1996). ``The Inconsistency of Common Scale Estimators When Output Prices Are Unobserved and Endogenous,'' \emph{Journal of Applied Econometrics}, 11(4), 343--361.

\bibitem{DLS2021} De Loecker, J.\ and C.~Syverson (2021). ``An Industrial Organization Perspective on Productivity,'' in \emph{Handbook of Industrial Organization}, Vol.~4, eds.\ K.~Ho, A.~Hortacsu, and A.~Lizzeri, Elsevier, 141--223.

\bibitem{LiuWangXu2022} Liu, L., Y.~Wang, and Y.~Xu (2022). ``A Practical Guide to Counterfactual Estimators for Causal Inference with Time-Series Cross-Sectional Data,'' \emph{American Journal of Political Science}, 66(1), 160--176.

\bibitem{OlleyPakes1996} Olley, G.S.\ and A.~Pakes (1996). ``The Dynamics of Productivity in the Telecommunications Equipment Industry,'' \emph{Econometrica}, 64(6), 1263--1297.

\bibitem{Oster2019} Oster, E.\ (2019). ``Unobservable Selection and Coefficient Stability: Theory and Evidence,'' \emph{Journal of Business \& Economic Statistics}, 37(2), 187--204.

\bibitem{Titl2023} Titl, V.\ (2023). ``The Effect of the Single-Bidding Ban on Procurement Prices,'' Working Paper.

\end{thebibliography}
\endgroup

\end{document}
